\chapter{الكسور والقوى}\label{ch:g9:fractions}

الكسور والقوى أدوات الحساب اليومية: لتقاسم
المقادير، ومقارنة النسب، وكتابة الأعداد الكبيرة جدًّا أو الصغيرة
جدًّا. ويرسّخ هذا الفصل قواعد الحساب بها ---
مع كتابة كل خطوة وسيطة --- ويقدّم \omterm{def:g9:fractions:scientific}{الترميز
العلمي}.

\section{الحساب بالكسور}

\begin{definition}[الكسر]\label{def:g9:fractions:fraction}
\emph{الكسر}\index{كسر} $\dfrac ab$ (مع $b \neq 0$) هو
\omterm{def:g3:division:remainder}{خارج قسمة} $a$ على $b$. ويتساوى كسران حين يُحصل أحدهما من
الآخر بضرب \omterm{def:g4:fractions:def}{البسط} \omterm{def:g4:fractions:def}{والمقام} (أو قسمتهما) في
العدد نفسه غير المعدوم:
\[
\frac ab = \frac{a \times k}{b \times k} \qquad (k \neq 0).
\]
\end{definition}

\begin{example}\label{ex:g9:fractions:simplify}
اِختزل $\dfrac{42}{56}$ خطوة خطوة:
\[
\frac{42}{56} = \frac{42 \div 2}{56 \div 2} = \frac{21}{28}
= \frac{21 \div 7}{28 \div 7} = \frac34 .
\]
(وفي \cref{ch:g9:arith} سيقوم القاسم المشترك الأكبر بهذا في خطوة
واحدة.)
\end{example}

\begin{method}[جمع الكسور وطرحها]\label{met:g9:fractions:add}
\begin{enumerate}
  \item ضع الكسرين على \emph{\omterm{def:g4:fractions:def}{مقام} مشترك} (وهو \omterm{def:g5:division:multiple}{مضاعف}
        مشترك للمقامين)؛
  \item واجمع البسطين أو اِطرحهما، محتفظًا \omterm{def:g4:fractions:def}{بالمقام}؛
  \item واِختزل النتيجة إن أمكن.
\end{enumerate}
\end{method}

\begin{example}\label{ex:g9:fractions:add}
اِحسب $\dfrac56 + \dfrac34$. \omterm{def:g4:fractions:def}{ومقام} مشترك للعددين $6$ و $4$ هو
$12$:
\[
\frac56 + \frac34
= \frac{5 \times 2}{6 \times 2} + \frac{3 \times 3}{4 \times 3}
= \frac{10}{12} + \frac{9}{12}
= \frac{19}{12}.
\]
واِحسب $2 - \dfrac37$. اُكتب $2$ كسرًا على \omterm{def:g4:fractions:def}{المقام} $7$:
\[
2 - \frac37 = \frac{14}{7} - \frac37 = \frac{11}{7}.
\]
\end{example}

\begin{omfigure}
\begin{tikzpicture}[scale=0.55]
  \foreach \i in {0,...,11}
    \draw[omExo] ({mod(\i,12)},0) rectangle ({mod(\i,12)+1},1);
  \foreach \i in {0,...,9}
    \fill[omDef!35] (\i,0) rectangle (\i+1,1);
  \draw[thick] (0,0) rectangle (12,1);
  \node[below, font=\small] at (6,-0.1)
    {$\frac56 = \frac{10}{12}$: 10 خانات من 12};
  \begin{scope}[yshift=-2.4cm]
  \foreach \i in {0,...,11}
    \draw[omExo] (\i,0) rectangle (\i+1,1);
  \foreach \i in {0,...,8}
    \fill[omThm!35] (\i,0) rectangle (\i+1,1);
  \draw[thick] (0,0) rectangle (12,1);
  \node[below, font=\small] at (6,-0.1)
    {$\frac34 = \frac{9}{12}$: 9 خانات من 12};
  \end{scope}
\end{tikzpicture}

{\small لماذا \omterm{def:g4:fractions:def}{المقام} المشترك: تقطيع الكلّين إلى $12$ خانة
متساوية يجعل الكسرين قابلين للمقارنة والجمع --- ومعًا،
$10 + 9 = 19$ جزءًا من اِثني عشر.}
\end{omfigure}

\begin{method}[ضرب الكسور وقسمتها]\label{met:g9:fractions:mult}
\begin{enumerate}
  \item للضرب، اِضرب البسطين معًا والمقامين
        معًا: $\dfrac ab \times \dfrac cd = \dfrac{ac}{bd}$
        (واِختزل \emph{قبل} الضرب كلما أمكن)؛
  \item وللقسمة، اِضرب في \omterm{def:g8:fractions:inverse}{مقلوب} القاسم:
        $\dfrac ab \div \dfrac cd = \dfrac ab \times \dfrac dc$.
\end{enumerate}
\end{method}

\begin{example}\label{ex:g9:fractions:mult}
\[
\frac{7}{15} \times \frac{25}{14}
= \frac{7 \times 25}{15 \times 14}
= \frac{7 \times 25}{15 \times 2 \times 7}
= \frac{25}{30} = \frac56,
\]
حيث اِختزلنا على العامل المشترك $7$، ثم على $5$. وإليك
قسمة:
\[
\frac{3}{4} \div \frac{9}{8}
= \frac34 \times \frac89
= \frac{3 \times 8}{4 \times 9}
= \frac{24}{36} = \frac23 .
\]
\end{example}

\section{القوى}

\begin{definition}[القوة]\label{def:g9:fractions:power}
من أجل عدد حقيقي $a$ وعدد صحيح موجب $n$، تكون
\emph{القوة}\index{قوة} $a^n$ \omterm{def:g2:mult:def}{جداء} $n$ عاملًا يساوي كل منها
$a$:
\[
a^n = \underbrace{a \times a \times \dots \times a}_{n \text{ عاملًا}} .
\]
وباصطلاح $a^0 = 1$ (من أجل $a \neq 0$)، وتدلّ \omterm{def:g8:powers:def}{الأُسُس} السالبة على
المقلوبات:
\[
a^{-n} = \frac{1}{a^n}.
\]
\end{definition}

\begin{example}
$2^4 = 2 \times 2 \times 2 \times 2 = 16$؛ \quad
و $(-3)^2 = 9$ لكن $-3^2 = -(3^2) = -9$ (\omterm{def:g8:powers:def}{فالأُسّ} يربط قبل
إشارة الناقص)؛ \quad و $10^{-3} = \frac{1}{1000} = 0.001$؛ \quad
و $5^1 = 5$ و $5^0 = 1$.
\end{example}

\begin{theorem}[قواعد الأُسُس]\label{thm:g9:fractions:rules}
من أجل كل $a$ و $b$ غير معدومين وكل عددين صحيحين $m$ و $n$:
\[
a^m \times a^n = a^{m+n},
\qquad
\frac{a^m}{a^n} = a^{m-n},
\qquad
\left(a^m\right)^n = a^{mn},
\qquad
(ab)^n = a^n b^n .
\]
\end{theorem}

\begin{proof}
من أجل \omterm{def:g8:powers:def}{الأُسُس} الموجبة، عُدّ العوامل. ففي القاعدة الأولى:
$a^m \times a^n$ هو $m$ عاملًا $a$ يتبعها $n$ عاملًا $a$، أي في المجموع
$m + n$ عاملًا. وفي الثالثة:
$\left(a^m\right)^n$ يكرّر $n$ مرة كتلةً من $m$ عاملًا، فيعطي
$mn$ عاملًا. والقاعدتان الأخريان مشابهتان؛ والامتداد إلى \omterm{def:g8:powers:def}{الأُسّ} المعدوم
\omterm{def:g8:powers:def}{والأُسُس} السالبة يُتحقَّق منه من $a^{-n} = \frac{1}{a^n}$ (وهو يجعل
الاصطلاح $a^0 = 1$ الخيار المتّسق الوحيد، لأن
$a^n \times a^0$ يجب أن يكون $a^{n+0} = a^n$).
\end{proof}

\begin{example}\label{ex:g9:fractions:rules}
بسّط خطوة خطوة:
\[
\frac{3^7 \times 3^{-2}}{3^4}
= \frac{3^{7 + (-2)}}{3^4}
= \frac{3^5}{3^4}
= 3^{5-4} = 3^1 = 3 .
\]
ومع حرفين: $(2a^3)^2 \times a = 4 a^6 \times a = 4a^7$.
\end{example}

\section{قوى العشرة والترميز العلمي}

\begin{proposition}[قوى العشرة]\label{prop:g9:fractions:ten}
من أجل عدد صحيح موجب $n$: يُكتب $10^n$ هكذا: «$1$ يتبعه $n$
صفرًا»، و $10^{-n} = 0.0\dots01$ حيث يقع $1$ في المرتبة العشرية
ذات الرتبة $n$. وضرب \omterm{def:g6:decimals:places}{عدد عشري} في $10^n$ (أو في $10^{-n}$) ينقل
نقطته العشرية $n$ مرتبة إلى اليمين (أو إلى اليسار).
\end{proposition}

\begin{definition}[الترميز العلمي]\label{def:g9:fractions:scientific}
\emph{الترميز العلمي}\index{ترميز علمي} \omterm{def:g6:decimals:places}{لعدد عشري} موجب
هو الطريقة الوحيدة لكتابته على الصورة
\[
a \times 10^n,
\qquad\text{حيث } 1 \leq a < 10 \text{ و } n \text{ عدد صحيح.}
\]
\end{definition}

\begin{example}\label{ex:g9:fractions:scientific}
بُعد الأرض عن الشمس نحو $149\,600\,000$ كم $= 1.496 \times
10^8$ كم. \omterm{def:g6:measure:volume}{وحجم} جزيء الماء نحو $0.000\,000\,000\,28$ م
$= 2.8 \times 10^{-10}$ م. وللحساب بأعداد كهذه، جمّع
الأجزاء العشرية وقوى العشرة:
\[
(3 \times 10^5) \times (2.5 \times 10^{-2})
= (3 \times 2.5) \times 10^{5 + (-2)}
= 7.5 \times 10^{3}.
\]
\end{example}

\begin{method}[وضع عدد في الترميز العلمي]\label{met:g9:fractions:scientific}
\begin{enumerate}
  \item اُنقل \omterm{def:g4:decimals:point}{النقطة العشرية} إلى ما بعد أول رقم غير معدوم مباشرةً؛ فيعطي
        هذا العامل $a$ حيث $1 \leq a < 10$؛
  \item وعُدّ عدد المراتب التي اِنتقلت إليها النقطة: فذلك العدد هو $n$،
        موجبًا إن كان العدد الأصلي $\geq 10$، وسالبًا إن كان
        $< 1$؛
  \item وتحقّق: فالمقدار $a \times 10^n$ يجب أن يعيد العدد الأصلي.
\end{enumerate}
\end{method}

\section{تمارين}

\begin{exercise}[$\star$]\label{exo:g9:fractions:1}
اِحسب وأعطِ النتيجة كسرًا مختزلًا تمامًا:
\[
\frac13 + \frac15, \qquad
\frac72 - \frac53, \qquad
\frac{5}{12} + \frac38, \qquad
3 - \frac45 .
\]
\end{exercise}

\begin{exercise}[$\star$]\label{exo:g9:fractions:2}
اِحسب واِختزل:
\[
\frac59 \times \frac{27}{35}, \qquad
\frac{8}{15} \div \frac{4}{5}, \qquad
\frac23 \times \frac34 \times \frac45 .
\]
\end{exercise}

\begin{exercise}[$\star$]\label{exo:g9:fractions:3}
اِحسب:
\[
2^5, \qquad (-2)^4, \qquad -2^4, \qquad 4^{-2}, \qquad
\left(\tfrac23\right)^3, \qquad 7^0 .
\]
\end{exercise}

\begin{exercise}[$\star$]\label{exo:g9:fractions:4}
اُكتب على صورة \omterm{def:g9:fractions:power}{قوة} واحدة:
\[
5^3 \times 5^6, \qquad
\frac{7^9}{7^5}, \qquad
\left(2^3\right)^4, \qquad
\frac{3^2 \times 3^7}{3^5}, \qquad
2^5 \times 4 .
\]
\end{exercise}

\begin{exercise}[$\star$]\label{exo:g9:fractions:5}
اُكتب \omterm{def:g9:fractions:scientific}{بالترميز العلمي}: $45\,000$؛ و $0.0072$؛ و $302.5$؛
و $0.000\,000\,9$؛ واِثني عشر مليونًا.
\end{exercise}

\begin{exercise}[$\star\star$]\label{exo:g9:fractions:6}
اِحسب، وأعطِ النتيجة \omterm{def:g9:fractions:scientific}{بالترميز العلمي}:
\[
(2 \times 10^7) \times (4.5 \times 10^{-3}),
\qquad
\frac{9 \times 10^5}{3 \times 10^{-2}},
\qquad
(5 \times 10^3)^2 .
\]
\end{exercise}

\begin{exercise}[$\star\star$]\label{exo:g9:fractions:7}
خزّان ممتلئ إلى $\frac35$. وبعد إضافة $30$ لترًا يصير ممتلئًا إلى $\frac34$.
فما \omterm{def:g2:measure:units}{سعة} الخزّان؟ (عبّر أولًا عن \omterm{def:g9:fractions:fraction}{الكسر} المضاف من
الخزّان.)
\end{exercise}

\begin{exercise}[$\star\star$]\label{exo:g9:fractions:8}
تأكل عالية $\frac14$ من كعكة، ثم يأكل بدر $\frac13$ مما يبقى،
ثم تأكل كريمة $\frac12$ مما ما زال باقيًا. فأيّ \omterm{def:g9:fractions:fraction}{كسر} من
الكعكة يبقى في النهاية؟ (اِحسب \omterm{def:g3:division:remainder}{الباقي} بعد كل خطوة.)
\end{exercise}

\begin{exercise}[$\star\star$]\label{exo:g9:fractions:9}
ينتقل الضوء $3 \times 10^8$ متر في الثانية. والشمس تبعد نحو
$1.5 \times 10^{11}$ م. فكم من الوقت يستغرق ضوء الشمس ليبلغنا؟
أعطِ الجواب بالثواني (ولا حاجة إلى \omterm{def:g9:fractions:scientific}{الترميز العلمي})، ثم بالدقائق
والثواني.
\end{exercise}

\begin{exercise}[$\star\star\star$]\label{exo:g9:fractions:10}
أيّهما أكبر، $2^{30}$ أم $3^{20}$؟ (إرشاد: اُكتب الاثنين قوّتين
أُسّهما $10$ مستعملًا $\left(a^m\right)^n = a^{mn}$، وقارن $2^3$ مع
$3^2$.)
\end{exercise}

\section{مسألة: الشفرة السرّية للأعداد العشرية الدورية}

\begin{problem}[{مسألة نهاية الأسبوع --- لكل كسر كتابة عشرية
تتوقف أو تتكرّر، وكل عدد عشري دوري كسر،
والعدد $0.999\ldots$ يساوي $1$ بالضبط}]
\label{pb:g9:fractions:1}
اِقسم $3$ على $11$ فتقع الأرقام في ترنيمة:
$0.272727\ldots$ إلى الأبد. واِقسم $1$ على $4$ فتتوقف الأرقام
فجأةً: $0.25$. وتبرهن هذه المسألة على أن هذين هما السلوكان
الوحيدان الممكنان \omterm{def:g9:fractions:fraction}{للكسر} --- ثم تعكس الآلة:
فمن أيّ \omterm{def:g6:decimals:places}{عدد عشري} دوري تعيد بناء \omterm{def:g9:fractions:fraction}{الكسر}
المختبئ وراءه. وتحسم في الطريق، حسمًا نهائيًّا،
أكثر التساويات جدلًا في رياضيات المدرسة.

\medskip
\noindent\textbf{الجزء الأول --- من \omterm{def:g9:fractions:fraction}{الكسر} إلى \omterm{def:g6:decimals:places}{العدد العشري}.}
\begin{enumerate}
  \item اِحسب بالقسمة المطوّلة الكتابات العشرية للكسور
        $\frac14$ و $\frac13$ و $\frac56$ و $\frac{3}{11}$.
        فأيّها يتوقف، وأيّها يتكرّر، وبأيّ دورة؟
  \item اِشرح لماذا يجب أن تتوقف الكتابة العشرية \emph{لأيّ} \omterm{def:g9:fractions:fraction}{كسر}
        $\frac ab$ أو تتكرّر في النهاية. (ففي القسمة
        المطوّلة على $b$، أيّ قيم تستطيع البواقي
        المتعاقبة أن تأخذها؟ وماذا يحدث لحظة ظهور باقٍ
        للمرة الثانية؟)
  \item اِدفع قسمة $1$ على $7$ بعيدًا بما يكفي لتجد
        الدورة. فما طولها، وكيف يقارَن ذلك الطول
        بالحدّ الذي وعد به السؤال 2؟
  \item وبعض الكسور تتوقف لأن مقامها يمكن أن
        يُنفَخ إلى \omterm{def:g9:fractions:power}{قوة} للعشرة: $\frac14 =
        \frac{25}{100}$، و $\frac{3}{8} = \frac{375}{1000}$.
        اِشرح المحكّ، ولماذا لا يستطيع أيّ مضروب فيه صحيح
        أن يحوّل $3$ أو $6$ أو $11$ إلى $10$ أو $100$
        أو $1000, \dots$ (فبأيّ رقم كان سينتهي \omterm{def:g2:mult:def}{الجداء}؟).
  \item ودون قسمة، صنّف $\frac{9}{40}$ و $\frac{5}{12}$ و
        $\frac{33}{50}$ إلى «يتوقف» و«يتكرّر»، ثم أعطِ
        الكتابة العشرية للاثنين اللذين يتوقفان.
\end{enumerate}

\medskip
\noindent\textbf{الجزء الثاني --- من \omterm{def:g6:decimals:places}{العدد العشري} الدوري رجوعًا إلى
\omterm{def:g9:fractions:fraction}{الكسر}.}
\begin{enumerate}[resume]
  \item حيلة الإزاحة. ليكن $x = 0.7777\ldots$ اِحسب $10x$،
        ثم $10x - x$، واِستنتج $x$ كسرًا. وتحقّق
        بالقسمة.
  \item ليكن $x = 0.727272\ldots$ فأيّ \omterm{def:g9:fractions:power}{قوة} للعشرة تزيح $x$ بدورة
        كاملة واحدة بالضبط؟ اِستعملها لتكتب $x$
        كسرًا مختزلًا إلى أبسط صورة.
  \item وحالة مختلطة: $x = 0.58333\ldots$ (فالرقم $3$ وحده
        يتكرّر). اِحسب $1000x - 100x$ واِستنتج $x$
        كسرًا مختزلًا إلى أبسط صورة.
  \item والمشهورة: طبّق حيلة الإزاحة على
        $x = 0.9999\ldots$ فأيّ \omterm{def:g9:fractions:fraction}{كسر} --- وأيّ \emph{عدد}
        --- تجد؟ ووفّق بين النتيجة و
        حدسك، متذكّرًا الجار الشبح في
        \cref{pb:g6:decimals:1} \omterm{def:g3:division:remainder}{والباقي} المتقلّص في
        \cref{pb:g6:fractions:1}: أهو $0.999\ldots$ «تحت
        $1$ بقليل»، أم اِسم آخر للعدد $1$؟
  \item حوّل $2.454545\ldots$ و $0.123123123\ldots$ إلى
        كسرين مختزلين إلى أبسط صورة.
\end{enumerate}

\medskip
\noindent\textbf{الجزء الثالث --- القاموس مكتملًا.}
\begin{enumerate}[resume]
  \item توحي الأسئلة 6--10 بقاموس: فدورة من $k$
        رقمًا تتكرّر من \omterm{def:g4:decimals:point}{النقطة العشرية} تساوي تلك الدورة
        على $k$ تسعة ($0.727272\ldots = \frac{72}{99}$). اِستعمل
        القاموس لتتوقّع $0.037037\ldots$ كسرًا،
        واِختزل --- ثم اِعجب: أيّ \omterm{def:g9:fractions:fraction}{كسر} بسيط بشكل
        مدهش يرنّم «$037$»؟ وتحقّق منه بالقسمة.
  \item واجمع القاموس مع إزاحة: اُكتب
        $0.0272727\ldots$ كسرًا مختزلًا إلى أبسط صورة.
  \item والعدد $0.101001000100001\ldots$ (رقم $1$، ثم رقم
        $0$ واحد، ثم $1$، ثم رقمان $0$، ثم $1$، ثم ثلاثة
        أرقام $0$، وهكذا) لا يتوقف أبدًا. فهل هو \omterm{def:g6:decimals:places}{عدد عشري} دوري؟
        وماذا يقول الجزء الأول عنه حينئذٍ --- أيمكن أن يكون
        كسرًا؟ (فقد اِلتقيت للتوّ أول عدد \emph{أصمّ}
        مبرهن على صممه؛ وأشهرها يُمسَك
        في \cref{ch:g9:sqrt}.)
  \item وكم رقمًا للعدد $2^{30}$ تقريبًا؟ اِستعمل
        $2^{10} = 1\,024 \approx 1.024 \times 10^3$ وقواعد
        \omterm{def:g8:powers:def}{الأُسُس} (\cref{thm:g9:fractions:rules}) لتكتب
        $2^{30}$ \omterm{def:g9:fractions:scientific}{بالترميز العلمي} (تقريبًا)، و
        اُستنتج. (قارن \cref{exo:g9:fractions:10}.)
  \item الخاتمة: اُكتب $0.20262026\ldots$ كسرًا، و
        اُذكر الدرس ذا الاتجاهين في المسألة كلها: أيّ
        كتابات عشرية تكون كسورًا، وأيّ كسور لها
        أيّ كتابات عشرية؟
\end{enumerate}
\end{problem}
